% !TEX program = xelatex
% !TeX TS-program = xelatex
% =============================================================================
%  全国大学生数学建模竞赛（CUMCM）LaTeX 模板
%
%  Overleaf（必须）：
%    1. 只上传本 paper/ 目录（main.tex 作为主文档）
%    2. Menu -> Compiler -> 选 XeLaTeX（不要用 pdfLaTeX）
%    3. 同目录 latexmkrc 会强制 xelatex
%  根因：Overleaf 是 Linux，没有 SimSun/SimHei；pdfLaTeX 也无法加载本模板。
%  字体：默认 TeX 自带 Fandol（Overleaf 可用）；若本机有宋体/黑体则自动改用。
%  编译：xelatex main.tex （需两遍以生成目录）
% =============================================================================
\documentclass[UTF8, fontset=fandol, 12pt, a4paper]{ctexart}

% ---------- 页面：A4，上下左右 2.5cm ----------
\usepackage[a4paper, top=2.5cm, bottom=2.5cm, left=2.5cm, right=2.5cm]{geometry}

% ---------- 数学 ----------
\usepackage{amsmath,amssymb}

% ---------- 字体（fontspec + ctex，xelatex 编译） ----------
% 先以 Fandol 保证 Overleaf / TeX Live 可编译；Windows 本机若有中易字体再覆盖。
\IfFontExistsTF{SimSun}{%
  \setCJKmainfont[AutoFakeBold]{SimSun}%
  \setCJKsansfont[AutoFakeBold]{SimHei}%
  \setCJKfamilyfont{zhsong}[AutoFakeBold]{SimSun}%
  \setCJKfamilyfont{zhhei}[AutoFakeBold]{SimHei}%
  \IfFontExistsTF{KaiTi}{\setCJKfamilyfont{zhkai}[AutoFakeBold]{KaiTi}}{}%
  \IfFontExistsTF{FangSong}{\setCJKmonofont{FangSong}}{}%
}{}
\IfFontExistsTF{Times New Roman}{%
  \setmainfont{Times New Roman}%
}{%
  \IfFontExistsTF{TeX Gyre Termes}{\setmainfont{TeX Gyre Termes}}{}%
}
\IfFontExistsTF{Courier New}{%
  \setmonofont{Courier New}%
}{%
  \IfFontExistsTF{TeX Gyre Cursor}{\setmonofont{TeX Gyre Cursor}}{}%
}

% ---------- 正文：12.05pt，首行缩进 2em，两端对齐，行距 0.72em ----------
% 12pt 字号近似 12.05pt；0.72em 行距 => 行高 = 1.72 * 字号，\linespread 因子 ≈ 1.43
\linespread{1.43}
\setlength{\parindent}{2em}

% ---------- 标题编号（\ctexset 配置） ----------
% 一级标题用中文数字“一、二、三、”，二级及以下用 1.1 / 1.1.1
\ctexset{
  section/number       = \chinese{section},
  subsection/number    = \arabic{section}.\arabic{subsection},
  subsubsection/number = \arabic{section}.\arabic{subsection}.\arabic{subsubsection},
}

% ---------- 标题格式（titlesec 配置） ----------
\usepackage{titlesec}
\titleformat{\section}
  {\centering\fontsize{17.3pt}{20.8pt}\heiti}
  {\chinese{section}、}{1em}{}
\titleformat{\subsection}
  {\fontsize{14.45pt}{17.4pt}\heiti}
  {\arabic{section}.\arabic{subsection}}{1em}{}
\titleformat{\subsubsection}
  {\fontsize{12.05pt}{14.5pt}\heiti}
  {\arabic{section}.\arabic{subsection}.\arabic{subsubsection}}{1em}{}
\titlespacing{\section}       {0pt}{1.25em}{0.82em}
\titlespacing{\subsection}    {0pt}{1.15em}{0.55em}
\titlespacing{\subsubsection} {0pt}{1.15em}{0.55em}

% ---------- 目录（tocloft 点线引导） ----------
\usepackage{tocloft}
\renewcommand{\cftsecfont}{\heiti}
\renewcommand{\cftsecpagefont}{\heiti}
\renewcommand{\cftsecleader}{\cftdotfill{\cftdotsep}}
\renewcommand{\cftsubsecleader}{\cftdotfill{\cftdotsep}}
\renewcommand{\cftsubsubsecleader}{\cftdotfill{\cftdotsep}}

% ---------- 表格、插图与代码 ----------
\usepackage{booktabs}
\usepackage{array}
\usepackage{graphicx}
\usepackage{float}
\graphicspath{{figures/}{./}{../figures/}}
\usepackage{xcolor}
\usepackage{listings}
\lstset{
  basicstyle=\ttfamily\small,
  backgroundcolor=\color{black!3},
  frame=single,
  framesep=6pt,
  rulecolor=\color{black!30},
  framerule=0.8pt,
  breaklines=true,
  showstringspaces=false,
  columns=fullflexible,
  keepspaces=true,
}

% ---------- 页码：底部居中 ----------
\pagestyle{plain}

% ---------- 自定义命令（对应 Typst 函数） ----------
% \papertitle{标题}：居中 17.3pt 加粗标题
\newcommand{\papertitle}[1]{%
  \thispagestyle{empty}%
  \begin{center}
    {\fontsize{17.3pt}{22pt}\bfseries #1}
  \end{center}
  \vspace{1em}
}
% \abstractcn{摘要正文}{关键词}：居中“摘要”标题（14pt 加粗）+ 正文 + “关键字：”加粗 + 关键词 + 换页
\newcommand{\abstractcn}[2]{%
  \thispagestyle{empty}%
  \begin{center}{\fontsize{14pt}{17pt}\bfseries 摘要}\end{center}
  #1\par
  \vspace{1em}
  {\heiti 关键字：}#2\par
  \newpage
}
% \tocpage：目录页，标题“目录”居中 17.3pt 黑体加粗，点线引导
\makeatletter
\newcommand{\tocpage}{%
  \setcounter{page}{1}%
  \begin{center}
    {\fontsize{17.3pt}{20.8pt}\heiti 目录}
  \end{center}
  \vspace{0.5em}
  \@starttoc{toc}%
  \newpage
}
\makeatother
% \referencescn：参考文献节，标题“参考文献”（不带编号但进目录）
\newcommand{\referencescn}{%
  \addcontentsline{toc}{section}{参考文献}%
  % 只列入已核实来源，不编造条目。

\begin{thebibliography}{99}
\bibitem{luikov1975} Luikov A V. Systems of differential equations of heat and mass transfer in capillary-porous bodies[J]. International Journal of Heat and Mass Transfer, 1975, 18(1): 1-14.
\bibitem{hussain2003} Hussain M M, Dincer I. Two-dimensional heat and moisture transfer analysis of a cylindrical moist object subjected to drying: A finite-difference approach[J]. International Journal of Heat and Mass Transfer, 2003, 46(21): 4033-4039.
\bibitem{kaya2007} Kaya A, Aydin O, Dincer I. Numerical modeling of forced-convection drying of cylindrical moist objects[J]. Numerical Heat Transfer, Part A, 2007, 51(9): 843-854.
\bibitem{dasilva2014} da Silva W P, e Silva C M D P S, Gama F J A. Estimation of thermo-physical properties of products with cylindrical shape during drying: The coupling between mass and heat[J]. Journal of Food Engineering, 2014, 141: 65-73.
\bibitem{adrover2019} Adrover A, Brasiello A, Ponso G. A moving boundary model for food isothermal drying and shrinkage: General setting[J]. Journal of Food Engineering, 2019, 244: 178-191.
\end{thebibliography}
%
}
% \appendixcn：附录节，标题“附录 A  核心代码”
\newcommand{\appendixcn}{%
  \section*{附录 A\quad 核心代码}%
  \addcontentsline{toc}{section}{附录 A\quad 核心代码}%
  \vspace{1.1em}

下面给出径向有限体积中，内部界面通量与表面 Robin 通量的核心计算，用于说明正文离散格式。完整求解还包括网格生成、BDF 时间积分与结果输出。

\begin{lstlisting}[language=Python]
import numpy as np

def harmonic(a, b):
    return 2 * a * b / (a + b) if a > 0 and b > 0 else 0.0

def moisture_rhs(C, Ca, r, W, Dfun, hm):
    N = len(C) - 1
    dC = np.zeros_like(C)
    for i in range(N):
        rm = 0.5 * (r[i] + r[i + 1])
        Dm = harmonic(Dfun(C[i]), Dfun(C[i + 1]))
        Flux = -rm * Dm * (C[i + 1] - C[i]) / (r[i + 1] - r[i])
        if i == 0:
            dC[0] -= Flux / W[0]
        else:
            dC[i] -= Flux / W[i]
        dC[i + 1] += Flux / W[i + 1]
    dC[N] -= r[N] * hm * (C[N] - Ca) / W[N]
    return dC
\end{lstlisting}
%
}
% \threelinetable{标题}{列定义}{表头}{表格内容}：三线表，标题在表格上方居中 10.5pt 黑体加粗
% 调用示例：\threelinetable{表 1 示例}{ccc}{A & B & C}{1 & 2 & 3 \\ 4 & 5 & 6}
\newcommand{\threelinetable}[4]{%
  \begin{center}
    {\fontsize{10.5pt}{12.6pt}\heiti #1}\par
    \vspace{0.6em}
    \begin{tabular}{#2}
      \toprule
      #3 \\
      \midrule
      #4 \\
      \bottomrule
    \end{tabular}
  \end{center}
}

\begin{document}

% 封面（不编号）
\papertitle{圆柱形中药材热风烘干过程的径向热湿传递建模}

% 摘要页（不编号）
\abstractcn{%
针对圆柱形中药材的热风烘干过程，建立以半径为自变量的瞬态传热与水分扩散模型，用节点型有限体积法离散空间、用隐式BDF推进时间，按题目规定输出内部温度场与干基含水率场。问题一在预热平衡阶段采用附录给定的常物性，温度与水分分开积分；烘房温湿过程用保形三次Hermite插值（PCHIP）进入边界。1800\,s时中心与表面温度分别为33.5758\,$^\circ$C与36.7879\,$^\circ$C，失水集中在外侧，表面含水率降至1.5102\,kg/kg，中心在四位小数下仍为2.5500\,kg/kg。问题二从同一初值出发，全程改用随含水率、温度变化的经验物性，并对两场作联合非线性求解。3\,h时径向温差已很小，中心与表面含水率分别降至1.7662\,kg/kg与1.0081\,kg/kg。问题三沿用问题二的耦合方程，以全域最大含水率首次低于0.15\,kg/kg定义烘干结束时刻，并对4\,h后的恒温环境作延续假设。问题四根据给定的半径轨迹引入均匀比例收缩，在材料坐标下更换附录物性后重新计算达标时间。模型可复现题目表格，同时把蒸发潜热、端面效应和界面水分势等简化作为明确适用限制。%
}{%
热风烘干 \quad 热湿传递 \quad 有限体积 \quad 变物性 \quad 收缩}

% 目录页（页码从 1 开始）
\tocpage

\section{问题重述}
\subsection{问题背景}

干燥是决定中药材成品品质的关键工序之一。热风烘干通常包含预热平衡与恒温干燥两个阶段：通过调控烘房的温度和湿度，使圆柱形药材内部的热量向内传递、水分向外迁移，最终把各处干基含水率降到工艺要求以下。工艺参数选择不当会导致干燥效率低、能耗高和成品品质不稳定；完全依赖试验优化则成本高、周期长。因此需要建立能够给出药材内部温度与水分浓度时空分布的数学模型，用数值仿真得到干燥规律。

本题研究对象为长度 $L=25\,\mathrm{cm}$、半径 $R=2\,\mathrm{cm}$ 的近似圆柱药材。烘干开始时药材温度均匀为 $28^\circ\mathrm{C}$，干基含水率均匀为 $2.55\,\mathrm{kg/kg}$。烘房温度与水分浓度随时间的变化由附件给出；第四问另给出烘干过程中药材半径随时间的变化。相关热物性、扩散系数及对流换热、对流传质系数由各问附录给出。

\subsection{问题内容}

\textbf{问题一：}在预热平衡阶段，采用附录给出的常物性参数，建立药材温度 $T(r,t)$ 与干基含水率 $C(r,t)$ 的变化规律。按题目表格给出 $100,300,600,900,1200,1500,1800\,\mathrm{s}$ 时、到中心距离 $0,0.5,1,1.5,2\,\mathrm{cm}$ 处的温度和含水率，并保存 $1800\,\mathrm{s}$ 内每隔 $1\,\mathrm{s}$、半径每隔 $0.1\,\mathrm{cm}$ 的完整结果。

\textbf{问题二：}烘干过程一般持续 $2$--$3$ 天，预热与恒温阶段的物性不同。为简化问题，相关经验公式统一采用另一组随含水率和温度变化的公式，建立整个烘干过程前 $3\,\mathrm{h}$ 的温度与含水率模型。按表格给出每隔 $0.5\,\mathrm{h}$、上述五个半径处的结果，并保存每隔 $1\,\mathrm{s}$、半径每隔 $0.1\,\mathrm{cm}$ 的完整结果。

\textbf{问题三：}按烘干要求，药材各处水分浓度应低于 $0.15\,\mathrm{kg/kg}$。确定烘干所需时间，并按表格给出每隔 $6\,\mathrm{h}$ 及结束时刻、五个半径处的含水率。

\textbf{问题四：}实际烘干中药材因失水发生尺寸变化。根据给定的半径--时间数据，更换相应经验物性后重新确定烘干时长，并按动态半径输出含水率表格。


\section{问题分析}
本题不是平均失水曲线拟合，也不是烘房温度优化。待求量是圆柱内部的温度场与干基含水率场。毛细多孔体干燥可用温度与湿度的耦合传递来描述\cite{luikov1975}；本题附录已给出物性与对流系数，因此以题给公式为主模型，不把文献中的全部交叉项搬入。附件提供的是外部环境和几何轨迹，没有药材内部观测，因此模型有效性依靠守恒结构、解析退化核验、网格与时间加密以及物理界限检查，而不能声称已经得到实验标定。

药材长度与直径之比 $L/(2R)=6.25$。两端面积与侧面积之比为 $R/L=0.08$，预热阶段热扩散特征长度小于中截面到端面的距离，因此第一问以中部横截面的径向一维模型作为主线。该近似在数小时乃至数日尺度上会变弱，后文将其作为明确局限而不是已经证明可以忽略的事实。

\subsection{问题一的分析}

第一问时间仅为 $1800\,\mathrm{s}$，热物性取常数，水分扩散系数只依赖于含水率。常物性热方程存在 Bessel--Duhamel 解析核验；变扩散系数的水分问题一般没有同形式闭式解。表面换热与传质均为第三类边界，表面温度和含水率都是待求量，不能直接赋为烘房值。空气含湿量与药材干基含水率即使同以 $\mathrm{kg/kg}$ 表示，质量基准也不一定相同，本文将其作为题给等效传质势使用。

\subsection{问题二的分析}

第二问要求覆盖整个烘干过程的前 $3\,\mathrm{h}$，并改用随含水率变化的密度、比热、导热系数以及同时依赖含水率和温度的扩散系数。预热结束时刻的 $1800\,\mathrm{s}$ 只是第一问的输出期限，并不等于烘房已经进入 $50^\circ\mathrm{C}$ 恒温，因此第二问从 $t=0$ 用新物性重新积分，不拼接第一问数值场。两场通过物性相互耦合，必须联合求解。

\subsection{问题三的分析}

第三问在同一组变物性方程上继续积分，直到药材各处含水率低于 $0.15\,\mathrm{kg/kg}$。``各处''应取计算域内最大值，不能用截面平均或表面值代替。附件环境数据只到 $4\,\mathrm{h}$，之后需要明确的恒温延续假设。后期扩散系数随含水率下降而显著减小，不能用初始扩散率或前 $3\,\mathrm{h}$ 斜率外推总时长。

\subsection{问题四的分析}

第四问同时更换物性并引入半径收缩。附件给出的是外半径轨迹，内部变形不能唯一识别。本文采用均匀比例径向收缩，并引入材料坐标 $\xi=r/R(t)$，使计算域固定在单位区间。干基含水率不是体积浓度，收缩本身不会自动把 $C$ 放大；输出必须按当前物理半径取值，域外单元格不得填零。


\section{模型假设}
为在题给数据下得到可闭合、可复现的计算模型，本文采用以下假设。它们同时构成模型的适用范围。

\begin{enumerate}
  \item 药材可视为均质、各向同性的直圆柱，输出对应长度中点横截面；忽略轴向端部换热与传质。
  \item 第一至三问几何固定，$R=0.02\,\mathrm{m}$、$L=0.25\,\mathrm{m}$；第四问长度仍取定值，仅半径按给定轨迹作均匀比例收缩。
  \item 初始温度与干基含水率在药材内均匀，分别为 $28^\circ\mathrm{C}$ 与 $2.55\,\mathrm{kg/kg}$。
  \item 烘房温度与水分浓度由附件观测点经 PCHIP 插值进入边界，精确通过全部观测点，不对曲线作平滑或删点。
  \item 表面换热、传质采用牛顿冷却型第三类边界，系数取题给常数；空气含湿量与药材含水率之差作为等效传质驱动力，不另引入未给出的吸附等温线。
  \item 主计算模型不显式计入蒸发潜热、辐射和体积热源。表面失水计入质量平衡，但汽化耗热不计入能量方程。
  \item 水分扩散采用题给经验公式。第一问 $D=D(C)$；第二、三问 $D=D(C,T)$；第四问更换另一组 $D(C,T)$。温度进入指数时使用局部开尔文温度。
  \item 附件给出的 $\rho(C)$ 在能量方程中作为有效热密度使用，不把它同时解释为可变干物质密度。
  \item 第三、四问以计算域内最大含水率首次低于 $0.15\,\mathrm{kg/kg}$ 作为烘干结束判据；4\,h 后烘房取 $50^\circ\mathrm{C}$、$0.05\,\mathrm{kg/kg}$ 作为恒温延续，该延续不是未来实测。
\end{enumerate}


\section{符号说明}
文中主要符号见表~\ref{tab:symbols}。未列出的局部记号在首次出现时定义。

\begin{table}[htbp]
\centering
\caption{主要符号说明}
\label{tab:symbols}
\begin{tabular}{ccc}
\toprule
\textbf{符号} & \textbf{含义} & \textbf{单位} \\
\midrule
$r$ & 到圆柱轴线的径向距离 & $\mathrm{m}$ \\
$t$ & 时间 & $\mathrm{s}$ \\
$R,L$ & 药材半径、长度 & $\mathrm{m}$ \\
$T(r,t)$ & 药材温度 & $^\circ\mathrm{C}$ \\
$C(r,t)$ & 干基含水率 & $\mathrm{kg/kg}$ \\
$T_a(t),C_a(t)$ & 烘房温度、烘房水分浓度 & $^\circ\mathrm{C}$，$\mathrm{kg/kg}$ \\
$T_s,C_s$ & 表面温度、表面含水率 & $^\circ\mathrm{C}$，$\mathrm{kg/kg}$ \\
$\rho$ & 有效密度 & $\mathrm{kg/m^3}$ \\
$c_p$ & 有效比热容 & $\mathrm{J/(kg\cdot K)}$ \\
$k$ & 有效热传导系数 & $\mathrm{W/(m\cdot K)}$ \\
$D$ & 水分浓度扩散系数 & $\mathrm{m^2/s}$ \\
$h$ & 对流换热系数 & $\mathrm{W/(m^2\cdot K)}$ \\
$h_m$ & 对流传质系数 & $\mathrm{m/s}$ \\
$\alpha$ & 热扩散率 $k/(\rho c_p)$ & $\mathrm{m^2/s}$ \\
$\mathrm{Bi}$ & 毕渥数 $hR/k$ & --- \\
$\xi$ & 材料坐标 $r/R(t)$ & --- \\
$M(t)$ & 径向最大含水率 & $\mathrm{kg/kg}$ \\
$t_*$ & 最大含水率达到阈值的临界时刻 & $\mathrm{s}$ \\
\bottomrule
\end{tabular}
\end{table}


\section{问题一的模型建立与求解}
\subsection{模型建立}

预热阶段药材与热空气同时交换热量和水分。取中部横截面，用径向坐标 $r$ 与时间 $t$ 描述状态。计算区域为 $0\le r\le R=0.02\,\mathrm{m}$，$0\le t\le 1800\,\mathrm{s}$。附录给定
\[
\rho=820\,\mathrm{kg/m^3},\quad
c_p=2600\,\mathrm{J/(kg\cdot K)},\quad
k=0.36\,\mathrm{W/(m\cdot K)},
\]
\[
h=25\,\mathrm{W/(m^2\cdot K)},\qquad
h_m=8\times10^{-7}\,\mathrm{m/s},
\]
以及
\begin{equation}
D(C)=7\times10^{-9}\exp(-0.89/C)\quad(\mathrm{m^2/s}).
\label{eq:q1-D}
\end{equation}
热扩散率 $\alpha=k/(\rho c_p)=1.6886\times10^{-7}\,\mathrm{m^2/s}$。$1800\,\mathrm{s}$ 内热扩散特征长度约 $0.0174\,\mathrm{m}$，小于中截面到端面的距离 $L/2=0.125\,\mathrm{m}$，支持本问采用径向近似。

在相邻观测时刻之间，对烘房温度 $T_a(t)$ 与水分浓度 $C_a(t)$ 作 PCHIP 插值，使曲线精确通过全部观测点，并避免在单调段内产生新的过冲。

基准模型中，环层蓄热率等于径向导热净输入，水分质量浓度按固定参考干物质骨架书写。约去公共几何因子后得到
\begin{equation}
\rho c_p\frac{\partial T}{\partial t}
=\frac1r\frac{\partial}{\partial r}\left(rk\frac{\partial T}{\partial r}\right),
\qquad
\frac{\partial C}{\partial t}
=\frac1r\frac{\partial}{\partial r}\left[rD(C)\frac{\partial C}{\partial r}\right].
\label{eq:q1-pde}
\end{equation}
初值与中心对称条件为
\begin{equation}
T(r,0)=28,\quad C(r,0)=2.55,\qquad
\frac{\partial T}{\partial r}(0,t)=\frac{\partial C}{\partial r}(0,t)=0.
\label{eq:q1-ic}
\end{equation}
表面第三类边界为
\begin{equation}
-k\frac{\partial T}{\partial r}(R,t)=h\bigl[T(R,t)-T_a(t)\bigr],
\qquad
-D(C_s)\frac{\partial C}{\partial r}(R,t)=h_m\bigl[C_s(t)-C_a(t)\bigr].
\label{eq:q1-bc}
\end{equation}
在常物性、$D=D(C)$ 且不显式计入潜热时，温度场与含水率场相互独立，可分开积分。圆柱湿物料的径向扩散与表面交换已有同类研究\cite{hussain2003,dasilva2014}，本文只借鉴其守恒思路，不移植其他物料的经验参数。

\subsection{数值求解}

采用包含圆心与表面的节点型有限体积法。节点 $r_i=i\Delta r$，$i=0,\ldots,N$，$\Delta r=R/N$。控制体权重 $W_i=(b_i^2-a_i^2)/2$，其中 $a_i=\max(0,r_i-\Delta r/2)$，$b_i=\min(R,r_i+\Delta r/2)$。内部界面取调和平均扩散系数
\[
D_{i+1/2}=\frac{2D(C_i)D(C_{i+1})}{D(C_i)+D(C_{i+1})},
\]
向外通量为
\[
F^T_{i+1/2}=-r_{i+1/2}k\frac{T_{i+1}-T_i}{\Delta r},\qquad
F^C_{i+1/2}=-r_{i+1/2}D_{i+1/2}\frac{C_{i+1}-C_i}{\Delta r}.
\]
半离散后，中心节点满足
\[
\dot T_0=\frac{4\alpha}{\Delta r^2}(T_1-T_0),\qquad
\dot C_0=\frac{4D_{1/2}}{\Delta r^2}(C_1-C_0);
\]
表面节点把 Robin 通量写入右侧控制体。时间方向采用变阶变步长隐式 BDF，水分非线性用 Newton 迭代。选择隐式是因为显式稳定限制 $\Delta t\sim\Delta r^2/\alpha$ 远小于输出步长，而不是更换物理模型。

发表网格在规定的 $0.1\,\mathrm{cm}$ 输出点上加密表面层，节点数为 $2113$。相对容差 $10^{-10}$，温度绝对容差 $10^{-10}$，水分绝对容差 $10^{-12}$，最大时间步 $0.5\,\mathrm{s}$。

\subsection{解析核验}

常物性热方程在环境温度分段光滑时可用分离变量加 Duhamel 原理核验。令 $\mathrm{Bi}=hR/k$，特征根满足 $\mu J_1(\mu)=\mathrm{Bi}\,J_0(\mu)$。对从初值出发的单位阶跃响应
\begin{equation}
\theta(r,t)=1-\sum_{n=1}^{\infty}
\frac{2\mathrm{Bi}}{(\mu_n^2+\mathrm{Bi}^2)J_0(\mu_n)}
J_0\!\left(\mu_n\frac{r}{R}\right)
\exp\!\left(-\mu_n^2\frac{\alpha t}{R^2}\right).
\label{eq:bessel}
\end{equation}
本题 $T_a(0)=T_0$，故
\begin{equation}
T(r,t)=T_0+\int_0^t \theta(r,t-\tau)\,T_a'(\tau)\,\mathrm{d}\tau.
\label{eq:duhamel}
\end{equation}
$T_a'$ 由 PCHIP 导数给出，不对附件作平滑。水分只在把 $D$ 冻结为常数时有同形式级数；变 $D(C)$ 的正式答卷仍是数值解。

温度场与 Bessel--Duhamel 解在全部输出点上的最大绝对差为 $3.42\times10^{-7}\,{}^\circ\mathrm{C}$。相邻两套加密网格的全场最大差：温度 $5.90\times10^{-7}\,{}^\circ\mathrm{C}$，水分 $4.56\times10^{-6}\,\mathrm{kg/kg}$。题目要求的 $7\times5$ 表格在四位小数下不再变化。热量与归一化水分收支的相对残差分别为 $1.22\times10^{-7}$ 与 $2.10\times10^{-6}$。

\subsection{计算结果}

表~\ref{tab:q1T} 与表~\ref{tab:q1C} 给出四位小数结果。$1800\,\mathrm{s}$ 时烘房温度为 $41.513^\circ\mathrm{C}$，表面 $36.7879^\circ\mathrm{C}$，中心 $33.5758^\circ\mathrm{C}$。中心含水率的源值为 $2.549992$，四位显示仍为 $2.5500$；失水集中在外侧，表面降至 $1.5102\,\mathrm{kg/kg}$。完整时空场按题目要求写入结果文件。

\begin{table}[htbp]
\centering
\caption{30分钟内药材的温度（单位：$^\circ\mathrm{C}$）}
\label{tab:q1T}
\begin{tabular}{cccccc}
\toprule
时间/s & $0\,\mathrm{cm}$ & $0.5\,\mathrm{cm}$ & $1.0\,\mathrm{cm}$ & $1.5\,\mathrm{cm}$ & $2.0\,\mathrm{cm}$ \\
\midrule
100  & 28.0000 & 28.0003 & 28.0038 & 28.0319 & 28.1795 \\
300  & 28.0406 & 28.0634 & 28.1513 & 28.3673 & 28.8486 \\
600  & 28.4533 & 28.5361 & 28.8039 & 29.3155 & 30.1641 \\
900  & 29.3243 & 29.4583 & 29.8757 & 30.6164 & 31.7295 \\
1200 & 30.5429 & 30.7099 & 31.2225 & 32.1130 & 33.4295 \\
1500 & 31.9960 & 32.1871 & 32.7667 & 33.7476 & 35.1211 \\
1800 & 33.5758 & 33.7724 & 34.3642 & 35.3626 & 36.7879 \\
\bottomrule
\end{tabular}
\end{table}

\begin{table}[htbp]
\centering
\caption{30分钟内药材的水分浓度（单位：$\mathrm{kg/kg}$）}
\label{tab:q1C}
\begin{tabular}{cccccc}
\toprule
时间/s & $0\,\mathrm{cm}$ & $0.5\,\mathrm{cm}$ & $1.0\,\mathrm{cm}$ & $1.5\,\mathrm{cm}$ & $2.0\,\mathrm{cm}$ \\
\midrule
100  & 2.5500 & 2.5500 & 2.5500 & 2.5500 & 2.2470 \\
300  & 2.5500 & 2.5500 & 2.5500 & 2.5492 & 2.0508 \\
600  & 2.5500 & 2.5500 & 2.5500 & 2.5353 & 1.8770 \\
900  & 2.5500 & 2.5500 & 2.5497 & 2.5045 & 1.7547 \\
1200 & 2.5500 & 2.5500 & 2.5482 & 2.4646 & 1.6586 \\
1500 & 2.5500 & 2.5499 & 2.5445 & 2.4206 & 1.5787 \\
1800 & 2.5500 & 2.5497 & 2.5383 & 2.3755 & 1.5102 \\
\bottomrule
\end{tabular}
\end{table}

场上 $D(C)$ 的最小值约为 $3.883\times10^{-9}\,\mathrm{m^2/s}$，约为 $D(C_0)$ 的 $0.786$；表面处变 $D$ 与冻结 $D$ 的含水率差约为 $0.0419$。因此不能把冻结扩散系数的 Bessel 解当作本问水分答卷。内部接近初值处 $D$ 几乎不变，失水锋面由外向内推进。


\section{问题二的模型建立与求解}
\subsection{模型建立}

第二问要求描述整个烘干过程的前 $3\,\mathrm{h}$。附录改用随含水率变化的热物性，以及同时依赖含水率和温度的扩散系数：
\begin{equation}
\begin{aligned}
\rho&=650+128C,\\
c_p&=1450+2736\frac{C}{1+C},\\
k&=0.21+0.38\frac{C}{1+C},\\
D&=2.4\times10^{-3}\exp(-0.45/C)\exp\bigl[-3850/(T+273.15)\bigr].
\end{aligned}
\label{eq:q2-prop}
\end{equation}
其中 $T$ 以摄氏度存储，指数中的温度化为开尔文。对流系数仍取 $h=25\,\mathrm{W/(m^2\cdot K)}$、$h_m=8\times10^{-7}\,\mathrm{m/s}$。控制方程为
\begin{equation}
\rho(C)c_p(C)\frac{\partial T}{\partial t}
=\frac1r\frac{\partial}{\partial r}\left[rk(C)\frac{\partial T}{\partial r}\right],
\qquad
\frac{\partial C}{\partial t}
=\frac1r\frac{\partial}{\partial r}\left[rD(C,T)\frac{\partial C}{\partial r}\right].
\label{eq:q2-pde}
\end{equation}
初值仍为 $T=28^\circ\mathrm{C}$、$C=2.55$，中心对称，表面 Robin 边界与\eqref{eq:q1-bc}相同，但 $k$ 与 $D$ 取表面局部状态。$\rho(C)$ 只作有效热密度，不把 $\rho/(1+C)$ 当作可变干密度。

第一问 $1800\,\mathrm{s}$ 时烘房仅约 $41.5^\circ\mathrm{C}$，不是恒温阶段的分界点。因此本问从 $t=0$ 用\eqref{eq:q2-prop}重新积分到 $10800\,\mathrm{s}$，不拼接第一问场。两场通过 $\rho,c_p,k,D$ 双向耦合，对状态向量 $(T_0,C_0,\ldots,T_N,C_N)$ 作联合隐式 BDF。界面导热与扩散系数仍用调和平均。环境在 $[0,10800]\,\mathrm{s}$ 上对附件观测点作 PCHIP。

\subsection{计算结果}

发表网格节点数为 $1057$。与更粗网格比较，全场最大差为温度 $2.45\times10^{-6}\,{}^\circ\mathrm{C}$、水分 $1.49\times10^{-5}\,\mathrm{kg/kg}$，论文 $6\times5$ 表四位小数不变。制造解短时检验给出温度最大差 $6.1\times10^{-8}\,{}^\circ\mathrm{C}$、水分 $9.0\times10^{-9}$。水分收支相对残差 $4.3\times10^{-7}$，计入 $B'(C)\dot C$ 修正后的有效热收支相对残差 $3.6\times10^{-8}$。

表~\ref{tab:q2T} 与表~\ref{tab:q2C} 给出前 $3\,\mathrm{h}$ 结果。$3\,\mathrm{h}$ 时烘房 $50.195^\circ\mathrm{C}$，表面 $49.9645^\circ\mathrm{C}$，中心 $49.8494^\circ\mathrm{C}$，径向温差已经很小。中心含水率降至 $1.7662\,\mathrm{kg/kg}$，表面降至 $1.0081\,\mathrm{kg/kg}$。相对第一问 $1800\,\mathrm{s}$ 中心几乎未干：升温使 $D$ 增大到约 $2.22\,D(C_0,28^\circ\mathrm{C})$，三小时扩散尺度足以影响轴心。冻结 $D=D(C_0,28^\circ\mathrm{C})$ 时中心含水率与变 $D$ 相差约 $0.40$，故不能用冻结解析解作为水分答卷。

$3\,\mathrm{h}$ 热扩散长度约 $5\,\mathrm{cm}$，已大于第一问，一维径向对端部的忽略需要在评价中保留。本问未引入蒸发潜热，也未使用第四问的收缩轨迹。

\begin{table}[htbp]
\centering
\caption{3小时内药材的温度（单位：$^\circ\mathrm{C}$）}
\label{tab:q2T}
\begin{tabular}{cccccc}
\toprule
时间/h & $0\,\mathrm{cm}$ & $0.5\,\mathrm{cm}$ & $1.0\,\mathrm{cm}$ & $1.5\,\mathrm{cm}$ & $2.0\,\mathrm{cm}$ \\
\midrule
0.5 & 32.1896 & 32.3824 & 32.9660 & 33.9608 & 35.4132 \\
1.0 & 40.3822 & 40.5542 & 41.0609 & 41.8783 & 43.0000 \\
1.5 & 45.8471 & 45.9351 & 46.1932 & 46.6055 & 47.1397 \\
2.0 & 48.4503 & 48.4881 & 48.5985 & 48.7736 & 49.0033 \\
2.5 & 49.4670 & 49.4793 & 49.5136 & 49.5645 & 49.6617 \\
3.0 & 49.8494 & 49.8553 & 49.8746 & 49.9096 & 49.9645 \\
\bottomrule
\end{tabular}
\end{table}

\begin{table}[htbp]
\centering
\caption{3小时内药材的水分浓度（单位：$\mathrm{kg/kg}$）}
\label{tab:q2C}
\begin{tabular}{cccccc}
\toprule
时间/h & $0\,\mathrm{cm}$ & $0.5\,\mathrm{cm}$ & $1.0\,\mathrm{cm}$ & $1.5\,\mathrm{cm}$ & $2.0\,\mathrm{cm}$ \\
\midrule
0.5 & 2.5499 & 2.5489 & 2.5256 & 2.3257 & 1.6486 \\
1.0 & 2.5257 & 2.4948 & 2.3578 & 2.0230 & 1.4711 \\
1.5 & 2.3861 & 2.3256 & 2.1344 & 1.8020 & 1.3476 \\
2.0 & 2.1708 & 2.1084 & 1.9236 & 1.6259 & 1.2311 \\
2.5 & 1.9566 & 1.9006 & 1.7360 & 1.4720 & 1.1166 \\
3.0 & 1.7662 & 1.7165 & 1.5702 & 1.3333 & 1.0081 \\
\bottomrule
\end{tabular}
\end{table}


\section{问题三的模型建立与求解}
\subsection{模型与环境延续}

第三问沿用\eqref{eq:q2-pde}--\eqref{eq:q2-prop}，几何与初值不变，仍不显式计入潜热和收缩。目标是确定使药材各处含水率低于 $0.15\,\mathrm{kg/kg}$ 的烘干时间，而不是另建优化模型。

附件环境序列只覆盖 $0$--$4\,\mathrm{h}$。末一小时温度与水分浓度的均值分别为 $50.00^\circ\mathrm{C}$ 与 $0.0500\,\mathrm{kg/kg}$，与题面所述恒温干燥阶段一致。因此 $0$--$4\,\mathrm{h}$ 仍使用观测点的 PCHIP 曲线，$4\,\mathrm{h}$ 之后取
\[
T_a(t)=50^\circ\mathrm{C},\qquad C_a(t)=0.05\,\mathrm{kg/kg}.
\]
切换处允许环境输入出现小幅跳变，药材内部 $T,C$ 保持连续。该延续是有数据支持的工作假设，不是未来实测。

$50^\circ\mathrm{C}$ 时，附录扩散系数在 $C=2.55$ 约为 $1.35\times10^{-8}\,\mathrm{m^2/s}$，在 $C=0.15$ 约为 $8.00\times10^{-10}\,\mathrm{m^2/s}$，后者只有前者的约 $6\%$。后期内部变干后扩散显著变慢，不能用初始 $D$ 或前 $3\,\mathrm{h}$ 的下降斜率外推总时长。

\subsection{达标判据}

定义径向最大含水率
\begin{equation}
M(t)=\max_{0\le r\le R}C(r,t),\qquad
t_*:=\inf\{t\ge0:M(t)<0.15\}.
\label{eq:q3-event}
\end{equation}
``各处低于阈值''必须用 $M(t)$ 判断。表格中的五个半径以及结果文件中的 $21$ 个采样点只是输出位置，不能代替全域最大值。若解沿径向向外不增，则 $M(t)=C(0,t)$，但仍需在加密网格上检查最大位置。

连续向下穿越时，临界时刻满足 $M(t_*)=0.15$；题目要求的是严格小于。正式报告时刻 $t_{\mathrm{rep}}$ 取略晚于 $t_*$ 且经未舍入场复核满足 $M(t_{\mathrm{rep}})<0.15$ 的时刻，并检查数值误差，不能只把临界根四舍五入后宣称已经达标。题面中的 $2$--$3$ 天是工艺经验范围，不作为强制截断或拟合目标。

\subsection{求解与输出}

数值格式与问题二相同：表面加密的径向有限体积、物性随局部状态更新、联合 BDF。在阈值附近限制步长，用符号变化括号定位事件，避免把输出间隔 $60\,\mathrm{s}$ 误当成求解步长。温度场作为内部状态保留，即使题目只要求导出含水率。

表~\ref{tab:q3C} 按题目格式列出每隔 $6\,\mathrm{h}$ 及结束时刻的含水率。当前全文稿先给出模型、判据与环境假设；结束时刻的具体小时数及表中数值，待按上述方程完成全程积分后填入，不在此处估计或编造。

\begin{table}[htbp]
\centering
\caption{药材烘干过程的水分浓度（单位：$\mathrm{kg/kg}$，结束时刻待补）}
\label{tab:q3C}
\begin{tabular}{cccccc}
\toprule
时间/h & $0\,\mathrm{cm}$ & $0.5\,\mathrm{cm}$ & $1.0\,\mathrm{cm}$ & $1.5\,\mathrm{cm}$ & $2.0\,\mathrm{cm}$ \\
\midrule
6  & --- & --- & --- & --- & --- \\
12 & --- & --- & --- & --- & --- \\
18 & --- & --- & --- & --- & --- \\
$\cdots$ & --- & --- & --- & --- & --- \\
结束时刻 & --- & --- & --- & --- & --- \\
\bottomrule
\end{tabular}
\end{table}


\section{问题四的模型建立与求解}
\subsection{收缩描述}

第四问根据给定半径轨迹重新确定烘干时长，并更换另一组经验物性。半径记录覆盖 $0$--$72\,\mathrm{h}$，由 $2\,\mathrm{cm}$ 单调降至 $1.198\,\mathrm{cm}$。附件只给出外半径，不能唯一识别内部非均匀变形。本文采用均匀比例径向收缩，并取长度 $L$ 不变。半径函数 $R(t)$ 对观测点作 PCHIP，保证 $R>0$ 且不产生非物理反弹。

干基含水率 $C$ 是水与干物质的质量比，不是单位体积含水量。在固定长度、干物质不流失且初始干物质密度均匀的条件下，径向固体速度取
\[
v(r,t)=\frac{\dot R(t)}{R(t)}r.
\]
引入材料坐标 $\xi=r/R(t)$。均匀收缩使材料速度与网格速度一致，对流项抵消，水分方程化为固定单位区间上的变系数扩散：
\begin{equation}
\frac{\partial U}{\partial t}
=\frac{1}{R(t)^2\xi}\frac{\partial}{\partial\xi}
\left[\xi D(U,\Theta)\frac{\partial U}{\partial\xi}\right],
\qquad 0<\xi<1,
\label{eq:q4-C}
\end{equation}
其中 $U(\xi,t)=C(R(t)\xi,t)$。相应的有效能量方程为
\begin{equation}
\rho(U)c_p(U)\frac{\partial\Theta}{\partial t}
=\frac{1}{R(t)^2\xi}\frac{\partial}{\partial\xi}
\left[\xi k(U)\frac{\partial\Theta}{\partial\xi}\right].
\label{eq:q4-T}
\end{equation}
不能只移动边界而漏掉材料运动，也不能在已经抵消后再重复加入几何对流项\cite{adrover2019}。

\subsection{物性与边界}

附录给出
\begin{equation}
\begin{aligned}
\rho&=760+90U,\\
c_p&=1850+2150\frac{U}{1+U},\\
k&=0.12+0.20\frac{U}{1+U},\\
D&=4.2\times10^{-4}\exp(-0.30/U)\exp\bigl[-3850/(\Theta+273.15)\bigr].
\end{aligned}
\label{eq:q4-prop}
\end{equation}
初值 $U=2.55$、$\Theta=28^\circ\mathrm{C}$，中心导数为零。表面等效边界写到 $\xi=1$ 上：
\[
-\frac{k}{R}\frac{\partial\Theta}{\partial\xi}=h(\Theta_s-T_a),\qquad
-\frac{D}{R}\frac{\partial U}{\partial\xi}=h_m(U_s-C_a).
\]
环境与问题三相同：$0$--$4\,\mathrm{h}$ 用观测曲线，之后取 $50^\circ\mathrm{C}$、$0.05\,\mathrm{kg/kg}$。$\rho c_p$ 仍解释为有效热容量；若把同一 $\rho$ 同时当成真实湿密度，并与固定长度、均匀收缩和给定 $R(t)$ 联立，会在后期过约束，因此不以该解释作为严格多相守恒模型。

同一温度、含水率下，本问扩散系数与问题三所用公式之比约为 $0.175\exp(0.15/C)$，在湿态约为 $0.19$、在 $C=0.15$ 时约为 $0.48$。收缩使扩散距离缩短，但物性同时变慢，第三、四问时长之差不能全部归因于半径变化。

\subsection{达标与输出}

达标判据与\eqref{eq:q3-event}相同，最大值在 $\xi\in[0,1]$ 上取。输出按当前物理半径：$r=\xi R(t)$。表~\ref{tab:q4C} 在每 $6\,\mathrm{h}$ 及结束时刻给出 $0,0.5,\ldots\,\mathrm{cm}$ 且不超过当时半径的位置，并单列实际表面值；域外单元格留空，不得填 $0$ 或空气含湿量。例如 $6\,\mathrm{h}$ 时 $R=1.374\,\mathrm{cm}$，有效列仅为 $0$、$0.5$、$1.0\,\mathrm{cm}$ 及表面。

求解仍在固定 $\xi$ 网格上作有限体积与联合 BDF，每步更新 $R(t)$、$k$、$D$ 和 $\rho c_p$。退化检验包括：$R$ 为常数时回到固定域模型；关闭边界交换且 $C$ 均匀时，纯收缩不改变干基含水率。结束时刻的数值同样待积分完成后填入。

\begin{table}[htbp]
\centering
\caption{收缩条件下烘干过程的水分浓度（单位：$\mathrm{kg/kg}$，结束时刻待补）}
\label{tab:q4C}
\begin{tabular}{cccccc}
\toprule
时间/h & $0\,\mathrm{cm}$ & $0.5\,\mathrm{cm}$ & $1.0\,\mathrm{cm}$ & $\cdots$ & 药材表面 \\
\midrule
6  & --- & --- & --- & --- & --- \\
12 & --- & --- & --- & --- & --- \\
18 & --- & --- & --- & --- & --- \\
$\cdots$ & --- & --- & --- & --- & --- \\
结束时刻 & --- & --- & --- & --- & --- \\
\bottomrule
\end{tabular}
\end{table}


\section{敏感性分析}
\subsection{网格与时间步}

问题一比较三套表面加密网格。相邻两套的全场最大差已低于 $2\times10^{-5}$，题目 $7\times5$ 表的四位小数不再随再加密而改变。将最大时间步由 $0.5\,\mathrm{s}$ 收紧到 $0.25\,\mathrm{s}$ 后，温度与水分的最大变化分别为 $2.87\times10^{-8}\,{}^\circ\mathrm{C}$ 与 $1.18\times10^{-9}\,\mathrm{kg/kg}$，论文表四位不变。问题二的粗细网格比较给出同样结论：发表表对进一步加密不敏感。因此表~\ref{tab:q1T}--表~\ref{tab:q2C} 的四位显示可以作为当前模型的稳定输出。

\subsection{扩散系数非线性}

问题一若把 $D$ 冻结为 $D(C_0)$，表面含水率在 $1800\,\mathrm{s}$ 仍会与变 $D$ 相差约 $0.04$。问题二中升温使 $D$ 增大、失水使 $D$ 减小，两者竞争后中心含水率与冻结 $D$ 相差约 $0.40$。这说明变扩散系数不是可有可无的修正，正式水分场必须保留题给 $D(C)$ 或 $D(C,T)$。

\subsection{后期环境与一维近似}

第三、四问在 $4\,\mathrm{h}$ 后采用 $50^\circ\mathrm{C}$、$0.05\,\mathrm{kg/kg}$。末小时均值与该设定几乎重合，但把最后观测点原样延续、或改用末小时均值，都会使达标时间发生有界移动。正式补入表~\ref{tab:q3C}、表~\ref{tab:q4C} 时应并列给出这一情景差，避免把单一延续写成唯一物理事实。

一维径向模型忽略端面。第一问短时可以接受；问题二 $3\,\mathrm{h}$ 热扩散长度已约 $5\,\mathrm{cm}$，更长的干燥过程中端部影响会进入内部。若后续补算轴对称 $(r,z)$ 模型，应以同一物性和两端 Robin 条件比较最大含水率到达时间，而不是在没有对照的情况下宣称一维结果已经偏于保守。


\section{模型评价与推广}
\subsection{模型优点}

\begin{enumerate}
  \item 四问共用同一径向守恒框架，由常物性、变物性耦合到收缩变域逐步加机制，假设和方程可以逐条对应到题给附录，而不是更换一套互不衔接的算法名称。
  \item 空间用节点有限体积，边界通量有直接物理含义，便于检查热量与水分收支；时间用隐式 BDF，适合表面薄层引起的刚性。
  \item 第一问温度与 Bessel--Duhamel 解一致，冻结扩散系数问题可与级数解对照，从而把离散误差与物理简化分开讨论。
  \item 输出半径落在网格节点上，题目要求的四位小数表经过网格与时间加密后保持稳定。
\end{enumerate}

\subsection{模型缺点}

\begin{enumerate}
  \item 主模型不显式计入蒸发潜热。表面失水已经进入质量平衡，但汽化耗热未进入能量方程。若把给定密度解释为初始湿密度，则预热阶段隐含潜热与显热增量为同一量级，不能以``预热时间短''断言潜热可忽略\cite{dasilva2014}。题面缺少吸附等温线和相容的界面传质定义，目前无法在不引入额外参数的前提下把潜热闭合进主计算。
  \item 空气含湿量与药材干基含水率被当作同一传质势作差。同类圆柱干燥研究中的边界往往使用物料平衡含水率，而不是空气含湿量本身\cite{dasilva2014,kaya2007}。
  \item 一维径向近似忽略端面；长时干燥中该误差会积累。
  \item 第四问的均匀收缩与有效热密度是在仅有外半径轨迹时的闭合办法，不是已经由内部变形观测证实的力学模型。
\end{enumerate}

\subsection{推广}

同一套有限体积--BDF 框架可以推广到：在能量边界中加入与失水通量一致的潜热项；用轴对称二维模型估计端部影响；在获得吸附等温线后把 $C_a$ 换成平衡含水率；以及把已知 $R(t)$ 换成由局部含水率驱动的收缩律。推广时必须补充相应参数，并重新做守恒与加密检验，不能直接套用其他农产品的经验系数。


\newpage
\referencescn
\newpage
\appendixcn

\end{document}
